\documentclass[12pt]{article}

\usepackage[english]{babel}
\usepackage[margin=1in, top=0.7in]{geometry}
\usepackage{microtype}
\usepackage{parskip}
\usepackage{xcolor}
\usepackage[colorlinks=true,urlcolor=blue]{hyperref}

\definecolor{accentblue}{RGB}{31, 73, 125}

\begin{document}

% Header
\noindent
\begin{minipage}[t]{0.45\textwidth}
    {\LARGE \textbf{\textcolor{accentblue}{Elisabetta Roviera}}}
\end{minipage}
\hfill
\begin{minipage}[t]{0.45\textwidth}
    \raggedleft
    {\small
    Nichelino (TO), Italy \\
    +39 340 519 03 87 \\
    \href{mailto:elisabetta.roviera00@gmail.com}{elisabetta.roviera00@gmail.com}
    }
\end{minipage}

\vspace{0.3cm}
{\color{accentblue}\hrule}
\vspace{0.5cm}

\today

\bigskip

Dear Prof. Stein,

I am applying for the PhD fellowship in computational analysis and modelling of protein interaction. Having read about the Proteins in CONTEXT project and your group's recent work, I believe this is the research environment I have been looking for: one where mathematics, machine learning and biology genuinely converge — and where the goal is not to apply methods to data, but to build models that are honest about biological complexity and rigorous about their own limitations.

What interests me most is the challenge of modelling intrinsically disordered regions. Coming from a mathematical background, I find this particularly compelling: function emerging from flexibility rather than from a fixed structure is a setting where classical approaches reach their limits — and where, I think, the most interesting methodological questions are still open.

My background is in Mathematical Engineering (MSc, Politecnico di Torino, 110/110) — a degree combining rigorous mathematics with applied data science, spanning optimisation, statistical learning, and scientific computing. For my Master's thesis, supervised by Prof.\ Alfredo Benso in collaboration with clinicians at Rivoli Hospital, I developed a machine learning pipeline to classify normal versus cancer-adjacent breast tissue from high-dimensional DNA methylation data. Working across 866k CpG sites and multiple cohorts, I learned that doing this seriously meant engaging with the biology — understanding what batch effects mean clinically, and how to build a feature selection strategy that respects genomic structure, not just statistical variance. A central contribution was a Mixed-Integer Linear Programme compressing a 5,000-site panel to 30 biologically constrained CpGs. A paper from this work has been submitted to IWBBIO 2026. I am currently continuing this research as a Research Fellow under Prof.\ Benso.

This experience clarified something I had been working towards without fully naming it. I chose mathematics because it fascinated me, but I always wanted to use it for something larger — something that could matter to people. During my thesis I found that direction: working at the intersection of mathematics, machine learning and biology, on problems that are both technically hard and clinically meaningful. The more I worked on it, the more I realised how much I still want to learn — both computationally and on the biology side. A PhD at the boundary between these disciplines, in an environment that takes both seriously, is what I am looking for.

I am aware that my background is not in structural biology or molecular dynamics. Learning the biology properly is not a limitation I am trying to work around — it is part of why I want to do a PhD. I want to deepen my computational skills and finally build the biological foundation to go with them.

I would be glad to discuss my background and motivation further.

\bigskip

Yours sincerely,

\bigskip\bigskip

Elisabetta Roviera

\end{document}